\subsection{Inverses with Odd Powers} 
Since odd power functions are one-to-one and odd roots are their inverses, we can write the inverses of odd power functions in general.
\begin{Example}Let $f(x)=2(x-3)^3+6$. Find $f\inv$.
\begin{lpic}{}{7}
\foreach \y/\lbl in {1/Swap $x$ and $y$,
										2/Isolate the 3rd power,
										5/Rewrite using a radical function,
										6/Finish solving for $y$,
										7/Rewrite in function notation}
	\regtextleft{0.25}{-\y}{\lbl};
\end{lpic}
\end{Example}
\noi Since we know that an odd power function is also the inverse of an odd root, we can also write the inverse of an odd radical function with the same logic.
\begin{Example}Let $g(x)=\frac{1}{2}\sqrt[3]{5x+1}-1$. Find $g\inv$.
\begin{lpic}{}{9}
\foreach \y/\lbl in {1/Swap $x$ and $y$,
										2/Isolate the radical,
										5/Rewrite as a power function,
										6/Finish solving for $y$,
										9/Rewrite in function notation}
	\regtextleft{0.25}{-\y}{\lbl};
\end{lpic}
\end{Example}